\documentclass[a4paper,11pt]{ltjsarticle}
\usepackage{luatexja}
\usepackage{luatexja-fontspec}
\usepackage{amsmath,amssymb,amsthm}
\usepackage{mathtools}
\usepackage{booktabs}
\usepackage{longtable}
\usepackage{array}
\usepackage{hyperref}
\usepackage{graphicx}
\usepackage{float}
\usepackage{geometry}
\geometry{margin=25mm}

\hypersetup{
  colorlinks=true,
  linkcolor=blue,
  urlcolor=blue,
  citecolor=blue
}

\theoremstyle{definition}
\newtheorem{theorem}{定理}[section]
\newtheorem{proposition}[theorem]{命題}
\newtheorem{lemma}[theorem]{補題}
\newtheorem{corollary}[theorem]{系}
\newtheorem{definition}[theorem]{定義}
\newtheorem{remark}[theorem]{注意}
\newtheorem{observation}[theorem]{観察}

\setlength{\parindent}{1em}
\setlength{\parskip}{0.5em}

\providecommand{\tightlist}{\setlength{\itemsep}{0pt}\setlength{\parskip}{0pt}}
\providecommand{\real}[1]{#1}
\begin{document}

\section{ブラックホール熱力学の上位時空構造からの写像としての再解釈：仮説駆動型研究プログラム（論文1：俯瞰・仮説提示）}\label{ux30d6ux30e9ux30c3ux30afux30dbux30fcux30ebux71b1ux529bux5b66ux306eux4e0aux4f4dux6642ux7a7aux69cbux9020ux304bux3089ux306eux5199ux50cfux3068ux3057ux3066ux306eux518dux89e3ux91c8ux4eeeux8aacux99c6ux52d5ux578bux7814ux7a76ux30d7ux30edux30b0ux30e9ux30e0ux8ad6ux65871ux4fefux77b0ux4eeeux8aacux63d0ux793a}

\textbf{著者}：木原 範昭（WF System Co., Ltd.）
\textbf{ORCID}：\href{https://orcid.org/0009-0004-6753-4020}{0009-0004-6753-4020}
\textbf{日付}：2026年4月
\textbf{DOI}：\href{https://doi.org/10.5281/zenodo.19837588}{10.5281/zenodo.19837588}

\begin{center}\rule{0.5\linewidth}{0.5pt}\end{center}

\subsection{要旨}\label{ux8981ux65e8}

ブラックホール熱力学を2つの条件付き命題の視点から探索する仮説駆動型研究プログラムを提案する：有効4次元時空内のブラックホールは高次元幾何学的構造の射影像でありうる（仮説1）こと、時空が
Planck スケールで離散的なら Hawking
蒸発が高次元領域境界からの離散単位ドリフトに対応しうる（仮説2）こと。両仮説は分野で広く共有される3つの前提から従う：BH
エントロピーが幾何学的、蒸発が地平面に局在、時空が Planck
スケールで離散的でありうる。本第一論文はプログラムを既存文献に位置づけ、2仮説を精密に明確化し、5つの姉妹論文での検証構造を予告し、予期される批判に応答する。本論文はブラックホール熱力学を導出しない；研究プログラムを定義する。

\begin{center}\rule{0.5\linewidth}{0.5pt}\end{center}

\subsection{§1. 序論}\label{ux5e8fux8ad6}

ブラックホール熱力学の記述は、1970年代初頭の Bekenstein, Hawking,
Bardeen--Carter--Hawking
により確立されたが、古典・半古典重力におけるもっとも顕著で、かつ十分に理解されていない発展の一つとして残されている。外部から見たブラックホールは熱的天体として振る舞う：地平面の面積に比例するエントロピーを持ち、質量に反比例する特徴的な温度で輻射し、有限の寿命で蒸発する。これらの性質は固定された曲がった背景上の場の量子論の枠内で導出され、その微視的起源
--- Bekenstein--Hawking
エントロピーに数えられるべきマイクロ状態は何か、Planck
スケールで蒸発を引き起こす機構は何か ---
は、重力・量子論・情報を結ぶ問題として、今なお業界で交渉中である。

この問題への複数の研究プログラムが存在する。't Hooft (1993) と Susskind
(1995) のホログラフィー原理は Maldacena (1997) の AdS/CFT
対応で具体的に実現され、重力理論のバルク自由度はその境界に符号化されると提唱する。ループ量子重力、因果動力学的三角分割、因果集合理論はそれぞれ異なる方法で時空が
Planck
スケールで本質的に離散的であると仮定し、この離散性からブラックホール熱力学を導出することを目指す。高次元一般化
--- Tangherlini ブラックホール (1963) と Myers--Perry 回転解 (1986) ---
は4次元公式を任意時空次元に拡張し、ホログラフィーや弦理論的アイデアの標準的試験場となっている。これらは各々部分的洞察を提供するが、いずれも完全で観測的に検証可能な記述には至っていない。

本論文は補完的方向を探る仮説駆動型研究プログラムを提案する。我々は業界で広く共有されていると考える3つの前提から出発する：ブラックホールエントロピーが幾何学的であること（Bekenstein）、蒸発が地平面スケールの動力学を反映すること（Hawking）、時空が
Planck
スケールで離散的でありうること（離散量子重力プログラムの共通前提）。これらから2つの仮説を抽出する。第一は、有効な4次元時空内で見られたとき、ブラックホールは独立な4次元実体ではなく高次元幾何学的構造の射影像でありうるということ。第二は、時空が離散的最小単位から構成されるなら、これら単位の高次元構造内領域からのドリフトが、4次元射影で観測される
Hawking 蒸発に対応しうるということ。

本論文はブラックホール熱力学を導出しないことを冒頭で強調する。研究プログラムを定義する。数学的・物理的内容
--- 中心投影下の計量、4次元球内の単位立方体充填の組合せ論的構造、5次元
Tangherlini エントロピーとの漸近対応、曲率半径の自己整合的動的決定 ---
は5つの姉妹論文で展開され、各々自己完結的で各々単一の限定された主張を行う。本論文は接続組織を提供する：プログラムを既存文献内に位置づけ、2つの仮説を検証可能な精度で明確化し、全シーケンスにわたる検証構造を予告する。

3つの特徴がプログラムを区別する。第一に、ホログラフィー原理や AdS/CFT
対応を仮定しない。ここでの4+1次元構造は単一の素朴な構成 ---
\(\mathbb{R}^5\) に埋め込まれた半径 \(R\) の4次元球面 \(S^4\)
から4次元主観座標図への中心（グノモン）投影 --- から生じ、誘導計量の
Schwarzschild
様の振る舞いはこの構成の帰結であって入力ではない。第二に、プログラムの離散時空側は特定の量子重力の微視理論を前提としない。単位立方体が4次元球内に飛び出さずに詰め込まれる要件と、結果として生じる整数格子組合せ論のみを使用する。第三に、結果として生じる熱力学の自然な比較対象は4+1次元
Tangherlini ブラックホールであり、馴染みの 3+1 次元 Schwarzschild
ブラックホールではない。これは戦略的選択である：プログラムを高次元ブラックホールの形式機構が既に十分に発達した高次元重力文献（Randall--Sundrum,
ADD 模型, Myers--Perry）の中に位置づける。

本第一論文の貢献は控えめである。ブラックホール熱力学の未解決問題を解くと主張せず、新しい観測的シグネチャを予言せず、既存プログラムを置き換えない。仮説的枠組みを明示し、その性質として、もし2つの仮説が有用であると判明するなら、最小限の幾何学的・組合せ論的仮定から確立された
Tangherlini
結果との定量的対応がプログラムから生じる。枠組みが有用であるかは姉妹論文が答えようとする問いであり、本論文はその問いを枠付けるのみである。

論文構成。第2節で本プログラムが最終的に接触し回復する必要のある、確立されたブラックホール熱力学の結果をレビューする。第3節で2つの仮説を精密に述べる。第4節で検証プログラムを概説し、どの論文がどの仮説にどのレベルの厳密さで貢献するかを特定する。第5節でプログラムが直面しうる主要批判を予期し先回りで応答する。第6節で結論する。

\subsection{§2.
既存ブラックホール熱力学のレビュー}\label{ux65e2ux5b58ux30d6ux30e9ux30c3ux30afux30dbux30fcux30ebux71b1ux529bux5b66ux306eux30ecux30d3ux30e5ux30fc}

本節では、検証プログラムが最終的に回復または接触する必要のある、確立されたブラックホール熱力学の結果をレビューする。レビューは意図的に簡潔・中立的とする。歴史的発展の主要な系統に沿って構成し、§3
で導入する仮説に関連する結果を強調する。

\subsubsection{§2.1 Schwarzschild
計量と古典的記述}\label{schwarzschild-ux8a08ux91cfux3068ux53e4ux5178ux7684ux8a18ux8ff0}

Einstein 場方程式の最も単純な非自明真空解は 1916 年に導出された
Schwarzschild 解である。標準 \((t, r, \theta, \phi)\) 座標で
\[ds^2 = -f(r)\, dt^2 + f(r)^{-1}\, dr^2 + r^2 \, d\Omega_2^2, \qquad f(r) = 1 - \frac{2GM}{r}.\]
\(f(r)\) は \(r_s = 2GM\) で消え、Schwarzschild
地平面を定義する。古典的な試験粒子力学では地平面は特別な役割を持たない（測地線は適切な座標フレームで滑らかに横断する）が、以下でレビューする熱力学的・量子力学的記述では中心的重要性を獲得する。

\subsubsection{§2.2 Hawking
輻射：熱的蒸発}\label{hawking-ux8f3bux5c04ux71b1ux7684ux84b8ux767a}

1975 年、Hawking は曲がった時空上の場の量子論の形式機構を Schwarzschild
背景に適用し、際立った結果を導出した：Schwarzschild ブラックホールは温度
\[T_H = \frac{\hbar \kappa}{2\pi k_B c} = \frac{\hbar c^3}{8\pi G M k_B}\]
で熱輻射を放射する。ここで \(\kappa\) は地平面での表面重力。輻射は
Planck
的スペクトルを持ち、ブラックホールからエネルギーを運び去り、有限の蒸発寿命
\(\tau \sim G^2 M^3 / (\hbar c^4)\) をもたらす。

温度は質量に反比例：小さいブラックホールほど熱く、速く蒸発する。蒸発の終点
--- Planck 質量ブラックホールの運命 --- は未解決問題のままである。

\subsubsection{§2.3 Bekenstein
エントロピー}\label{bekenstein-ux30a8ux30f3ux30c8ux30edux30d4ux30fc}

1973 年、Bekenstein
は熱力学的根拠から、ブラックホールが地平面面積に比例するエントロピーを持つに違いないと論じた。彼の提案を
Hawking の温度と組み合わせると Bekenstein--Hawking エントロピー
\[S_{BH} = \frac{k_B c^3 A}{4 G \hbar} = \frac{A}{4 \ell_P^2}\]
が得られる。ここで \(A = 4\pi r_s^2\) は Schwarzschild
地平面の面積、\(\ell_P = (G\hbar/c^3)^{1/2}\) は Planck
長。エントロピーは質量の 2 乗にスケールする。

\(S_{BH}\) の微視的起源 --- ブラックホールマイクロ状態として何を数えるか
---
は分野の主要な未解決問題。要求されるマイクロ状態数は膨大：太陽質量ブラックホールで
\(S_{BH}/k_B \sim 10^{77}\)。

\subsubsection{§2.4
ブラックホール熱力学の四法則}\label{ux30d6ux30e9ux30c3ux30afux30dbux30fcux30ebux71b1ux529bux5b66ux306eux56dbux6cd5ux5247}

Bardeen, Carter, Hawking (1973) は完全な熱力学的類比を確立：

\begin{itemize}
\tightlist
\item
  \textbf{第零法則}。表面重力 \(\kappa\)
  は定常地平面上で一定（平衡における温度均一性の類比）。
\item
  \textbf{第一法則}。\(dM = (\kappa / 8\pi G)\, dA + \Omega\, dJ + \Phi\, dQ\)（\(dU = T\, dS - p\, dV\)
  の類比）。
\item
  \textbf{第二法則}。地平面面積は古典過程で非減少；量子過程では一般化エントロピー
  \(S_{\mathrm{gen}} = S_{BH} + S_{\mathrm{matter}}\) が非減少。
\item
  \textbf{第三法則}。\(\kappa = 0\)
  の地平面（極限）には任意の有限物理過程で到達できない。
\end{itemize}

これらの法則は現在、古典・半古典の設定で頑健な熱力学的構造として理解されている。

\subsubsection{§2.5
ホログラフィー原理}\label{ux30dbux30edux30b0ux30e9ux30d5ux30a3ux30fcux539fux7406}

't Hooft (1993) と Susskind (1995)
はホログラフィー原理を提案：空間領域内の独立な量子力学的自由度の数は、Planck
単位の境界面積で抑えられる（体積ではない）。原理は Bekenstein--Hawking
エントロピーから動機づけられる：ブラックホールエントロピーが地平面面積で支配されるなら、任意領域内に収まる最大エントロピーも面積で抑えられる。

原理は、\(D\) 時空次元の重力物理の基本的記述が \((D-1)\)
次元曲面上の自由度を持つことを示唆する。それは下記の AdS/CFT
対応の概念的根源だが、論理的にはより一般。

\subsubsection{§2.6 AdS/CFT 対応}\label{adscft-ux5bfeux5fdc}

Maldacena (1997) は、\(\mathrm{AdS}_5 \times S^5\) 上のタイプ IIB
弦理論が、\(\mathrm{AdS}_5\) 境界上の \(\mathcal{N} = 4\) 超 Yang--Mills
理論と等価であると提案した。対応はホログラフィー原理の具体的実現：バルクの重力物理が境界ゲージ理論に符号化される。

ブラックホール熱力学の文脈では、AdS/CFT は明示的なマイクロ状態計数（BPS
セクター）と熱的 CFT
状態によるブラックホール地平面の双対記述を与える。対応はホログラフィー原理の最も成功した現代の実装。

\subsubsection{§2.7
高次元ブラックホール}\label{ux9ad8ux6b21ux5143ux30d6ux30e9ux30c3ux30afux30dbux30fcux30eb}

Tangherlini (1963) は Schwarzschild 解を任意時空次元 \(D\)
に一般化した。\(D = d + 1\) 次元での Schwarzschild--Tangherlini 計量：
\[ds^2 = -f_T(r)\, dt^2 + f_T(r)^{-1}\, dr^2 + r^2\, d\Omega_{d-1}^2, \qquad f_T(r) = 1 - \frac{\mu}{r^{d-2}}.\]
\(\mu\) は質量に比例。地平面は \(r_h = \mu^{1/(d-2)}\)、地平面の表面積は
\(S^{d-1}(r_h)\) の体積、Bekenstein--Hawking エントロピー公式は
\(S_{BH} = A_{d-1}/(4\ell_P^{d-1})\) に一般化。

我々の目的に関連するのは \(D = 5\)（\(d = 4\)）の場合：
\[f_T(r) = 1 - \frac{8GM}{3\pi r^2}, \qquad r_h^2 = \frac{8GM}{3\pi}, \qquad T_H = \frac{1}{2\pi r_h}, \qquad S_{BH} = \frac{2\pi^2 r_h^3}{4 G_5}.\]
これがここで展開するプログラムの自然な比較対象。

Myers と Perry (1986) は解を回転を含むよう一般化；Myers--Perry 解は Kerr
の高次元類似。

\subsubsection{§2.8
量子重力プログラム}\label{ux91cfux5b50ux91cdux529bux30d7ux30edux30b0ux30e9ux30e0}

時空を Planck スケールで本質的に離散と扱う複数のプログラム。

\emph{ループ量子重力}（Rovelli, Smolin 他、1990 年代）は正準量子重力の
Hilbert 空間を構築し、面積・体積演算子が Planck
スケールの最小単位を持つ離散スペクトルを持つ。Bekenstein--Hawking
エントロピーは地平面のスピンネットワーク状態の数え上げで回復される。

\emph{因果動力学的三角分割}（Ambjørn, Jurkiewicz, Loll、1998
年以降）は離散因果構造への非摂動経路積分を実装し、連続シミュレーションで創発的な4次元時空を回復する。

\emph{因果集合理論}（Bombelli, Lee, Meyer,
Sorkin、1987）は滑らかな時空多様体を、Planck
スケール事象間の因果関係を表す離散半順序で置き換える。

これらは §3.1 の前提 C（Planck
スケール離散性）を共有するが、具体的実装は大きく異なる。本論文の検証プログラムは特定実装にコミットせず、「離散性が関連スケールで意味を持つ」という抽象的前提のみを要求する。

\subsubsection{§2.9 総合}\label{ux7dcfux5408}

確立されたブラックホール熱力学の体系は以下の構造的特徴を持つ：

\begin{itemize}
\tightlist
\item
  幾何学的量に比例するエントロピー（前提 A）
\item
  起源が地平面に局在する半古典的蒸発機構（前提 B）
\item
  Planck スケール離散性を基本要素として援用する一連のプログラム（前提
  C）
\end{itemize}

これら3前提は共通の地点；プログラム間の差異は各々の具体的実装に関する。本論文の検証プログラムはこれら共通前提を抽出し、特定の幾何学的・組合せ論的実装
--- 高次元構造からの中心投影と整数格子の単位立方体充填 ---
を追求する。これは我々の知る限り文献では未追求である。

\subsection{§3. 2つの仮説}\label{ux3064ux306eux4eeeux8aac}

検証プログラムが組織される2つの仮説を述べる。両方とも条件付き命題として定式化する：現段階ではどちらも物理的根拠で真であると主張しない。

\subsubsection{§3.1 前提}\label{ux524dux63d0}

ブラックホール熱力学と量子重力の現代文献の相当部分で共有されていると考える3つの前提から出発する：

\textbf{前提
A}（\emph{エントロピーの幾何学的起源}）。ブラックホールのエントロピーは地平面に関連する幾何学的量で決定される。標準的定式化ではこれは
Planck
単位の地平面面積であるが、現段階では正確な幾何学的量を未指定のままにする。

\textbf{前提 B}（\emph{蒸発の地平面スケール起源}）。Hawking
輻射は地平面スケールでの動力学を反映し、無限遠ではない。蒸発を担う微視的機構が何であれ、それは地平面近傍領域に局在する。

\textbf{前提 C}（\emph{Planck スケール離散性}）。時空は Planck
スケールで実効的に離散でありうる。離散性の形は指定しない：十分微細な解像度で幾何学的量が整数値の組合せ論的データで記述可能であるという要請のみを行う。

前提 A は Bekenstein--Hawking 公式の合意的読解。前提 B は Hawking 1975
導出の操作的内容：関連する場の量子論計算は地平面近傍領域で行われ、結果として無限遠での熱スペクトルがそれら近傍モードの赤方偏移を反映する。前提
C
はループ量子重力、因果動力学的三角分割、因果集合理論の共通出発点：特定の実装にコミットせずに採用する。

\subsubsection{§3.2
仮説1：射影仮説}\label{ux4eeeux8aac1ux5c04ux5f71ux4eeeux8aac}

\textbf{仮説
1}。\emph{有効な4次元時空内で現れるブラックホールは、より高次元の時空に埋め込まれた幾何学的構造の射影像である。ブラックホールに関連する熱力学量
--- エントロピー、温度、蒸発レート ---
は自律的な4次元天体の性質ではなく、射影の性質である。}

仮説は前提 A
と整合する。4次元地平面が高次元幾何学的天体の射影なら、4次元での「地平面面積」はその高次元天体の何らかの幾何学的量に対応し、Bekenstein
様関係は4次元面積と高次元幾何学的量の関係として自然に解釈される。

仮説は高次元構造を\textbf{指定しない}。検証プログラムでは特定の具体的実現を扱う：\(\mathbb{R}^5\)
に埋め込まれた半径 \(R\) の4次元球面 \(S^4 \subset \mathbb{R}^5\)
から4次元主観座標図へのグノモン投影。これは可能な一つの実現であり、命題としての仮説はより一般。

仮説はホログラフィー原理と\textbf{等価でない}。ホログラフィーは境界をより低次元に置く；射影仮説はバルクをより\textbf{高次元}に置き、4次元物理を射影像と見なす。両者は論理的に独立：射影なしのホログラフィー、ホログラフィーなしの射影、両方、いずれもなし、を考えうる。

\subsubsection{§3.3
仮説2：離散ドリフト仮説}\label{ux4eeeux8aac2ux96e2ux6563ux30c9ux30eaux30d5ux30c8ux4eeeux8aac}

\textbf{仮説
2}。\emph{時空が離散最小単位から構成され、ブラックホールが有限高次元領域の射影なら、Hawking
蒸発はその領域からの離散単位のドリフトに対応する。具体的には、高次元領域の連続体積と、領域内に収まる離散単位の整数個数の差が、4次元射影で観測されるエントロピーと蒸発動力学を決定する。}

これは前提 B, C と整合。前提 B
は動力学が地平面に局在することを要求：離散ドリフト仮説は関連動力学を高次元領域の境界に局在させる。前提
C は Planck
スケール離散性を要求：仮説はこの離散性を真剣に取り、それが生じる計数問題を問う。

仮説はファズボールやマイクロ状態計数プログラムのいずれとも\textbf{等価でない}。ブラックホールが弦やブレーンから構成されると仮定しない；境界マイクロ状態を数えるために
AdS/CFT
を発動しない。「離散単位」は整数格子内の単位立方体と取られ、計数問題は4次元充填問題で、古典的かつ組合せ論的に
well-defined。それと熱力学の接続がプログラムの内容。

\subsubsection{§3.4 仮説の地位}\label{ux4eeeux8aacux306eux5730ux4f4d}

両仮説は条件付き命題として述べられる。高次元構造や離散単位の物理的実在性については主張しない；そのような構造を\textbf{仮定したとき}、確立された熱力学量を導出可能か、どの水準の一致でそうかを問うのみ。

姉妹論文では答えが部分的に肯定的であることを示す：

\begin{itemize}
\tightlist
\item
  中心投影下の誘導計量の Schwarzschild
  様の振る舞い（論文2）は計量構造のレベルで仮説1を支持。
\item
  漸近的充填解析（論文3）は体積不足 \(\Delta(R) \sim c\,R^3\) を
  \(c = 8/(3\pi)\) で与え、整合させるべき具体的量を提供。
\item
  4+1次元 Tangherlini Bekenstein--Hawking
  エントロピーとの対応（論文4）は仮説2の定性的スケーリングを定量的に実現。
\item
  動力学的自己整合解析（論文5）は標準 Tangherlini スケーリング
  \(R \propto M^{1/2}\), \(T_H \propto 1/\sqrt{M}\), 寿命
  \(\propto M^2\)
  を回復し、蒸発レート指数に1つの定量的不一致を残し、これを未解決問題として明示。
\end{itemize}

検証は部分的、不一致は明示。プログラムはブラックホール熱力学の問題を解いたとは主張しない。中心投影と離散性という最小限の仮定とともに2仮説を取ったとき、確立された結果に非自明に接触する計算が生まれる、と主張するのみ。

\subsection{§4.
検証プログラム：6篇の概観}\label{ux691cux8a3cux30d7ux30edux30b0ux30e9ux30e06ux7bc7ux306eux6982ux89b3}

本節では検証プログラムを概説し、どの論文がどの仮説にどの厳密さで貢献するかを特定する。各姉妹論文は自己完結的で、他論文と独立に読める。ここでの集合構造は記述的：プログラムを組織するが、個別論文の論理的前提条件ではない。

\subsubsection{§4.1
寄与のマッピング}\label{ux5bc4ux4e0eux306eux30deux30c3ux30d4ux30f3ux30b0}

§3 の2仮説への寄与をプログラム全体で以下のように扱う。

\textbf{仮説1（射影仮説）}：

{\def\LTcaptype{none} % do not increment counter
\begin{longtable}[]{@{}
  >{\raggedright\arraybackslash}p{(\linewidth - 4\tabcolsep) * \real{0.3333}}
  >{\raggedright\arraybackslash}p{(\linewidth - 4\tabcolsep) * \real{0.3333}}
  >{\raggedright\arraybackslash}p{(\linewidth - 4\tabcolsep) * \real{0.3333}}@{}}
\toprule\noalign{}
\begin{minipage}[b]{\linewidth}\raggedright
論文
\end{minipage} & \begin{minipage}[b]{\linewidth}\raggedright
寄与の種類
\end{minipage} & \begin{minipage}[b]{\linewidth}\raggedright
厳密さの水準
\end{minipage} \\
\midrule\noalign{}
\endhead
\bottomrule\noalign{}
\endlastfoot
2 & 数学的基盤：グノモン投影計量、Schwarzschild 類似形 &
厳密、定理-証明 \\
4 & 定量的対応：\(\Delta(R)\) vs.~4+1次元 Tangherlini エントロピー &
定量的、明示的数値整合 \\
5 & 動力学的整合：\(R\) の質量スケーリング & 定量的、1指数未解決 \\
6 & 統合または未解決問題 & （性格未定） \\
\end{longtable}
}

\textbf{仮説2（離散ドリフト仮説）}：

{\def\LTcaptype{none} % do not increment counter
\begin{longtable}[]{@{}lll@{}}
\toprule\noalign{}
論文 & 寄与の種類 & 厳密さの水準 \\
\midrule\noalign{}
\endhead
\bottomrule\noalign{}
\endlastfoot
3 & 数学的基盤：格子充填、漸近定数 \(c = 8/(3\pi)\) & 厳密、定理-証明 \\
4 & 同定：\(\Delta(R)\) をエントロピー量の候補として & 推測的 \\
5 & 動力学的役割：曲率半径進化におけるドリフト機構 & 発見的 \\
6 & 統合または未解決問題 & （性格未定） \\
\end{longtable}
}

\subsubsection{§4.2
論理的依存関係}\label{ux8ad6ux7406ux7684ux4f9dux5b58ux95a2ux4fc2}

姉妹論文は独立に読めるよう設計される。論文2
は中心投影下の計量構造のみ；論文3 は純粋に組合せ論的；論文4
は次元解析と特定の数値対応で両者を接続；論文5 は動力学を導入。系列

\[\text{論文 2, 3} \quad\rightarrow\quad \text{論文 4} \quad\rightarrow\quad \text{論文 5} \quad\rightarrow\quad \text{論文 6}\]

は自然な読書順だが、各論文は前者なしでも理解可能。

\subsubsection{§4.3
各仮説の妥当性領域}\label{ux5404ux4eeeux8aacux306eux59a5ux5f53ux6027ux9818ux57df}

\textbf{仮説1（射影）}：計量レベルで支持される --- 中心投影構成は
Schwarzschild
様の振る舞いと地平面構造を持つ計量を生む。これは特定の幾何学的構成についての数学的命題で、それ自体ではこの構成が物理的ブラックホールを記述すると確立しない。

\textbf{仮説2（離散ドリフト）}：スケーリングレベルで支持される ---
体積不足 \(\Delta(R)\) は 4+1次元 Tangherlini エントロピーと整合する
\(R^3\) スケーリングを持つ。\(\Delta(R)\)
を「離散ドリフトエントロピー」と解釈することは、数学的量を物理的にどう理解すべきかの仮説。論文5
で導出される動力学的スケーリング法則の整合性により部分的に支持されるが、離散単位が境界を超えてドリフトする正確な機構は導出されない。

\subsubsection{§4.4
蒸発レート指数}\label{ux84b8ux767aux30ecux30fcux30c8ux6307ux6570}

プログラムの主要な残存不一致は蒸発レート指数である。最簡単な動力学的仮定
--- Hawking 温度での地平面面積からの Stefan--Boltzmann 様輻射 --- は
\[\frac{dM}{dt} \propto -M^{n}, \qquad n = -1/2 \text{ または } 3/2 \text{（規約による）}\]
を与えるが、標準半古典 Hawking 計算は \(n = -1\)
を与える。不一致は特定の指数の整数オフセットで、符号誤りや \(\pi\)
の欠落ではない。

この不一致をプログラムの未解決問題として明示する。論文5
では可能な解決策を議論：

\begin{itemize}
\tightlist
\item
  輻射の幾何学への反作用のより注意深い扱い
\item
  最簡単なスケーリング論を超えた離散ドリフト機構の修正
\item
  射影された4次元理論における「蒸発」の意味の再解釈
\end{itemize}

プログラムは現段階で優先解決策を主張しない。論文6
は解決策を提案するか、不一致を精密に同定された未解決問題として残すかのいずれか。

\subsubsection{§4.5
プログラムが確立することと確立しないこと}\label{ux30d7ux30edux30b0ux30e9ux30e0ux304cux78baux7acbux3059ux308bux3053ux3068ux3068ux78baux7acbux3057ux306aux3044ux3053ux3068}

検証プログラム終了時に以下が確立される：

\begin{itemize}
\tightlist
\item
  中心投影構成の数学的整合性（論文2）
\item
  体積不足の漸近定数 \(c = 8/(3\pi)\) の正確な値（論文3）
\item
  \(\Delta(R)\) と 4+1次元 Bekenstein--Hawking エントロピーの \(R^3\)
  スケーリング対応（論文4）
\item
  Tangherlini スケーリング法則 \(R \propto M^{1/2}\),
  \(T_H \propto 1/\sqrt{M}\), 寿命 \(\propto M^2\) の回復（論文5）
\end{itemize}

以下は確立\textbf{されない}：

\begin{itemize}
\tightlist
\item
  離散ドリフト機構からの Bekenstein--Hawking
  エントロピーの第一原理的導出（スケーリングと主要係数のみ整合）
\item
  標準 Hawking
  計算との完全な一致（蒸発レート指数は未解決問題として残る）
\item
  ホログラフィー原理、AdS/CFT、標準量子重力プログラムとの関係の解決
\item
  物理的 3+1 次元ブラックホールの記述（自然な対象は 4+1 次元）
\end{itemize}

プログラムは、確立された Tangherlini
結果に対する\textbf{特定の幾何学的・組合せ論的枠組みの部分的検証}として理解されるべき。成功は枠組みが内的に整合し、さらなる発展に値することを示し；残存不一致は現在の定式化での枠組みの不足を示す。

\subsection{§5.
予期される批判への先回り}\label{ux4e88ux671fux3055ux308cux308bux6279ux5224ux3078ux306eux5148ux56deux308a}

プログラムが直面しうる主要批判に対応する。注意深い読者が提起しうる形で各批判を述べ、プログラムの応答を提供する。

\subsubsection{§5.1
なぜ4+1次元定式化か？}\label{ux306aux305c41ux6b21ux5143ux5b9aux5f0fux5316ux304b}

\textbf{批判}。観測される宇宙は3空間次元・1時間次元。詳細に研究され熱力学的性質が文献に固定されているブラックホールは
3+1 次元。なぜ本プログラムは自然な比較対象として 4+1 次元 Tangherlini
ブラックホールを取るのか？

\textbf{応答}。4+1 次元定式化は論文2 で採用される幾何学的構成 ---
\(\mathbb{R}^5\) に埋め込まれた4次元球面 \(S^4\)
から4次元主観座標図へのグノモン投影 ---
の帰結である。4次元主観座標図上の誘導計量は、独立な幾何学的パラメータとしての曲率半径
\(R\) を伴うとき、自然に5次元構造 \((x, y, z, t, R)\)
を生む。この構成で生じるブラックホール計量は次元的根拠で Tangherlini
類似であり、Schwarzschild 形ではない。

物理的時空次元が 4+1 であると主張しない。\textbf{論文2
の幾何学的構成が与えられたとき}、自然な比較は 4+1 ケースで、Tangherlini
熱力学との一致は構成自体の非自明な検証となると主張する。観測される 3+1
次元物理学との関係 --- Randall-Sundrum 様ブレーン局在化、ADD
様の大きな余剰次元、その他の機構 ---
は将来研究に委ね、本プログラム射程外。

この位置取りは高次元重力文献の相当部分と共有：Tangherlini, Myers-Perry,
Randall-Sundrum, ADD 全てが、3+1
次元での物理的実現は内的整合性とは別問題である一般化を研究している。

\subsubsection{§5.2
なぜホログラフィー原理を仮定しないか？}\label{ux306aux305cux30dbux30edux30b0ux30e9ux30d5ux30a3ux30fcux539fux7406ux3092ux4eeeux5b9aux3057ux306aux3044ux304b}

\textbf{批判}。ホログラフィー原理と AdS/CFT
はブラックホール熱力学に多大な進展をもたらした：詳細なエントロピー計数と明示的動力学的記述を含む。なぜ本プログラムはこれらを利用しないのか？

\textbf{応答}。プログラムは独立な研究経路として構築される。出発仮定 ---
高次元時空からの中心投影と Planck スケール離散性 ---
はホログラフィー原理と論理的に独立。結果として得られる描像がホログラフィーと整合的か、支持するか、矛盾するか、直交するかは、独立経路を完成させた後でのみ検討可能な問い。出発時点でホログラフィーを仮定するとその問いを覆い隠す。

ホログラフィー原理の不成立に依存するものは何もない。ホログラフィーが正しいなら、論文4・5
で導出される Tangherlini
熱力学との一致は何らかの機構でホログラフィー的計数と整合的であるはず：その正確な対応は本プログラムが解こうとしない未解決問題。

\subsubsection{§5.3
なぜ完全な定量的導出を含まないか？}\label{ux306aux305cux5b8cux5168ux306aux5b9aux91cfux7684ux5c0eux51faux3092ux542bux307eux306aux3044ux304b}

\textbf{批判}。論文1
は仮説提示のみ；定量的導出は姉妹論文へ延期。姉妹論文自体も蒸発レート指数（最簡単な扱いで
\(n = 3/2\)、標準 Hawking 計算で
\(n = -1\)）を一致させずに残す。プログラムは概略以上を提供するか？

\textbf{応答}。Yes。プログラムは以下を提供する：(i)
自身として厳密解析を許す精密な数学的枠組み ---
中心投影構成と整数格子充填問題；(ii) 体積不足の明示的漸近定数
\(c = 8/(3\pi)\) --- 解析的に導出され数値で確認；(iii) \(R^3\)
スケーリングおよび主要数値係数の 4+1 次元 Bekenstein--Hawking
エントロピーとの整合；(iv) 標準 Tangherlini スケーリング法則
\(R \propto M^{1/2}\), \(T_H \propto 1/\sqrt{M}\), 寿命 \(\propto M^2\)
の回復。

未解決の蒸発レート指数は真の未解決問題。プログラムは最簡単な動力学的仮定から得られる指数（\(n = 3/2\)）が標準半古典結果（\(n = -1\)）と一致しないことを明示し、より注意深い反作用または離散ドリフト機構の扱いが解決のために必要と予想する。この不一致への誠実さは方法論の一部であり、欠陥ではない。

\subsubsection{§5.4
プログラムはどのように検証可能か？}\label{ux30d7ux30edux30b0ux30e9ux30e0ux306fux3069ux306eux3088ux3046ux306bux691cux8a3cux53efux80fdux304b}

\textbf{批判}。仮説は条件付き命題として述べられ、プログラムは新しい観測を予言しない。どの意味で科学的か？

\textbf{応答}。プログラムは3つの異なる意味で検証可能：

\begin{enumerate}
\def\labelenumi{\arabic{enumi}.}
\item
  \emph{数学的整合性}。論文2, 3 は完全に数学的：中心投影計量、漸近定数
  \(c\)、\(N(k)\) の Jacobi
  様閉形式表現の任意の誤りは基礎的主張を反証する。数値計算は定義から再現可能；解析的導出は行ごとに検証可能。
\item
  \emph{定量的対応}。論文4, 5 は具体的数値関係 --- 定数
  \(c = 8/(3\pi)\)、比
  \(\Delta/S_{BH} \to 32/(3\pi) \approx 3.40\)、Tangherlini スケーリング
  \(R \propto M^{1/2}\) ---
  を生み、これらは整合するか否か。蒸発レート指数の例外を除き整合する。これら対応の不成立はプログラムの確立された熱力学との接触の主要主張を反証する。
\item
  \emph{未解決問題の特定}。論文6
  は残存不一致を形式化する。プログラムは蒸発レート指数を標準半古典計算との不一致の主要箇所として特定する。これ自体が有用な研究出力：完全整合のためにプログラムまたは標準計算のどちらに何を変更すべきかを指定する。
\end{enumerate}

新しい観測的シグネチャを予言しない。経験的物理学への寄与とは主張しない；特定の幾何学的仮定からの帰結の数学的・概念的探索と主張する。探索が経験的内容を持つかは将来研究の問いで、本プログラムの射程外。

\subsubsection{\texorpdfstring{§5.5
主張\textbf{しないこと}}{§5.5 主張しないこと}}\label{ux4e3bux5f35ux3057ux306aux3044ux3053ux3068}

誤読を防ぐため、プログラムが\textbf{主張しないこと}を明示する：

\begin{itemize}
\tightlist
\item
  物理的時空が 4+1 次元であるとは主張しない。
\item
  中心投影構成が重力の唯一・正しい幾何学的記述であるとは主張しない。
\item
  Bekenstein--Hawking エントロピーを第一原理から導出したとは主張しない。
\item
  蒸発レート指数が \(3/2\) であるとは主張しない；標準 \(n = -1\)
  との不一致を未解決問題として特定する。
\item
  ホログラフィック、ループ、因果集合、弦理論的プログラムと競合する・置き換える・反証するとは主張しない。
\end{itemize}

主張するのは、2つの条件付き命題（仮説1,
2）の連言、それらから従う数学的構造、確立された Tangherlini
結果との部分的定量的対応である。

\subsection{§6. 結論}\label{ux7d50ux8ad6}

ブラックホール熱力学を2つの条件付き命題の視点から探索する仮説駆動型研究プログラムを提案した。第一は、有効4次元時空内で見られるブラックホールが高次元幾何学的構造の射影像でありうること。第二は、時空が
Planck スケールで離散的なら、Hawking
蒸発が高次元領域境界からの離散単位ドリフトに対応しうること。

両仮説は、ブラックホール熱力学と量子重力の現代研究で広く共有される3つの前提から従う：BH
エントロピーが幾何学的、蒸発が地平面に局在、時空が Planck
スケールで離散的でありうる。両仮説は確立された結果を置き換えも矛盾もしない。我々の知る限り既存文献では追求されていない特定の探求経路
--- 高次元構造からの中心投影と整数格子組合せ論 --- を指定する。

検証プログラムは5つの姉妹論文から成り、各々自己完結的で各々単一の限定された主張を行う。プログラム終了までに以下が確立される：中心投影構成の整合性（論文2）、体積不足の漸近定数
\(c = 8/(3\pi)\) の正確な値（論文3）、4+1次元 Tangherlini
エントロピーとの対応（論文4）、Tangherlini
スケーリング法則の回復（論文5）。主要な残存不一致は蒸発レート指数で、これを精密に同定された未解決問題として残す。

プログラムはブラックホール熱力学の未解決問題を解かない。最小限の幾何学的・組合せ論的仮定から、確立された
Tangherlini
結果との非自明な定量的対応が導出可能であることを確立する。この対応が物理的に何かを示すのか、単なる数学的好奇心に過ぎないかは、プログラムが答えようとしない問い。

プログラムの価値は（あるとすれば）自己完結性と未解決問題の精密性にある。論文2,
3 の数学的内容は記述された仮定から完全に導出可能。論文4, 5
の数値的対応は明示的で、整合するか否かのいずれか。残存不一致は特定の定量水準で同定される。これらの特徴のいずれもホログラフィー、AdS/CFT、特定の量子重力プログラムに依存しない。プログラムは自己の内的整合性と直接比較する確立された結果との一致で立つか倒れるか。

3つの注記で締めくくる。

第一に、プログラムの 4+1 次元（3+1 ではなく）への位置取りは、論文2
の幾何学的構成により動機づけられた意図的戦略的選択である。時空の物理的次元数について主張しない；採用する構成の自然な舞台が
4+1 だから 4+1 で作業する。

第二に、ホログラフィー入力の不在も意図的。プログラムはホログラフィーが正しいか否かに独立に有用となるよう設計される。ホログラフィーが正しいなら、ここで導出される
Tangherlini
結果との一致はホログラフィー的計数と整合可能でなければならない；その整合はプログラム射程外。

第三に、方法論 --- 前提の明示、帰結の厳密な導出、不一致の精密な同定 ---
はプログラムの仮説を受け入れるかどうかにかかわらず有用となるよう意図する。仮説を拒否する読者でも自己完結的な数学的・組合せ論的枠組みを得て、その出力自身が興味深い可能性がある。

姉妹論文でのプログラムの展開は、論文2 の幾何学的構成から始まる。

\end{document}
